\section{{\fontsize{12pt}{14pt}\selectfont CONCLUSÃO}}
\vspace{-0.3cm}
\novoBlocoTexto {
    Este trabalho partiu do problema de pesquisa apresentado na Introdução, formulado na seguinte QP: ``\textit{Como proporcionar ao docente acesso organizado às dúvidas dos estudantes, a partir de seu próprio material didático, de modo a subsidiar o aprimoramento de suas práticas de ensino?}''. Para investigá-la, definiu-se como objetivo geral desenvolver e avaliar um protótipo web de um assistente educacional baseado em LLM e RAG, em o RAG recupera trechos dos materiais didáticos de professores e a LLM, com base nesses trechos, classifica e organiza as dúvidas dos estudantes nos tópicos do plano de ensino, disponibilizando essas informações em um painel para acompanhamento dos conteúdos questionados pelos estudantes.
}
\novoBlocoTexto {
    Os objetivos específicos foram atendidos ao longo do trabalho. O primeiro, identificar na literatura as dificuldades relacionadas à manifestação, recepção e organização das dúvidas discentes, foi cumprido pela revisão da literatura apresentada na Introdução e no referencial teórico, que mostrou que muitas dúvidas não chegam ao professor, seja pela resistência dos estudantes em manifestá-las, seja pelo uso de outras fontes de ajuda, como as ferramentas de inteligência artificial generativa, e que essas dúvidas são subsídios relevantes ao planejamento das aulas. O segundo, o levantamento de requisitos, resultou nos 14 requisitos funcionais e nos 9 requisitos não funcionais da seção~\ref{sec:levantamento-requisitos}. O terceiro, a modelagem das interfaces de usuário, foi atendido com os fluxos BPMN, que organizam as interações entre professor, aluno e sistema, e com as telas apresentadas na seção~\ref{sec:secao-5-1}. O quarto, o desenvolvimento do protótipo, foi realizado com o \textit{framework Next.js}, o \textit{Supabase}, a LLM \textit{Qwen2.5-7B-Instruct}, executada localmente, e o modelo de \textit{embeddings} \textit{multilingual-e5-base}; o protótipo registra as dúvidas, responde aos estudantes com base no material da disciplina e classifica cada dúvida em um dos tópicos do plano de ensino. O quinto foi atendido com a geração dos resumos individuais e com o painel docente organizado por tópicos, que exibe os resumos das perguntas e as perguntas originais sem vínculo com os estudantes. O Quadro~\ref{quad:atendimento-requisitos} mostra o atendimento de todos os requisitos.
}
\novoBlocoTexto {
    O desempenho da classificação foi avaliado em ambiente controlado, atendendo ao sexto objetivo específico deste trabalho. As 1.416 perguntas de estudantes de Estrutura de Dados do conjunto de dados de \textcite{zaman2024dataset}, traduzidas para o português, foram enviadas à plataforma pelo \textit{chat}, e o tópico atribuído pela LLM coincidiu com o tópico de referência em 83,33\% das perguntas, com precisão macro de 73,85\%, revocação macro de 84,76\% e macro F1-score de 74,90\%. O desempenho foi alto nos tópicos de conteúdo bem delimitado, como \textit{Hashing} (Espalhamento), Árvore binária de busca e Grafo, e menor nos tópicos que se sobrepõem, como Revisão de Python, Listas, Nós e Ponteiros. A análise dos erros apontou como causas prováveis das divergências a divisão diferente dos tópicos, os rótulos de referência que indicam o tema da aula, e não o assunto da pergunta, e os termos com mais de um sentido (seção~\ref{sec:analise-erros}); nos termos técnicos examinados, a tradução não se mostrou uma explicação relevante para as divergências. Com o protótipo desenvolvido e avaliado, o objetivo geral foi atendido.
}
\novoBlocoTexto {
    Em resposta à QP, os resultados indicam que é possível proporcionar ao docente acesso organizado às dúvidas dos estudantes, a partir de seu próprio material didático, por meio de um assistente em que o RAG recupera os trechos do material relacionados a cada dúvida e a LLM, com base neles, classifica a dúvida em um dos tópicos do plano de ensino e gera o seu resumo, apresentados ao professor em um painel organizado por tópico e sem a identificação dos estudantes. Essa organização dispensa a categorização manual das perguntas e, no experimento, apontou os mesmos seis tópicos com mais dúvidas da referência, ainda que em outra ordem. Como cerca de uma em cada seis perguntas foi apresentada em um tópico diferente do de referência, o painel deve ser entendido como subsídio à análise pedagógica, e não como diagnóstico definitivo, cabendo ao professor interpretar os dados e decidir as ações que podem aprimorar suas práticas de ensino.
}
\novoBlocoTexto {
    Esses resultados têm limitações: a avaliação usou perguntas traduzidas automaticamente, de uma única disciplina, enviadas por automação a partir de uma única conta, em uma única execução e com um único modelo, e um material elaborado a partir de apostilas para o experimento; além disso, a qualidade das respostas e dos resumos não foi avaliada (seção~\ref{sec:discussao-limitacoes}).
}
\novoBlocoTexto {
    Como trabalhos futuros, propõe-se o agrupamento, por similaridade, das dúvidas de um mesmo tópico, para que o professor visualize de forma reunida perguntas semelhantes feitas por diferentes estudantes, à semelhança do agrupamento por assuntos proposto por \textcite{angelo1993classroom} para as respostas da técnica \textit{muddiest point}. Esse recurso não foi incluído nesta versão porque o conjunto de dados utilizado na avaliação não dispõe de agrupamentos de referência. Também se propõem a inclusão de uma categoria para dúvidas não relacionadas à disciplina, que poderia ser avaliada com o arquivo de perguntas não relacionadas às disciplinas do mesmo conjunto de dados \parencite{zaman2024dataset}, a avaliação da qualidade das respostas e dos resumos gerados e a aplicação da plataforma em turmas reais, com perguntas em português e a participação de professores. Por fim, propõe-se a identificação da etapa em que surgem as divergências de classificação, a partir dos fragmentos recuperados registrados para cada pergunta.
}